\paragraph{Polars and orthogonals}

\begin{enumerate}
\def\labelenumi{\arabic{enumi})}
\setcounter{enumi}{4}
\tightlist
\item
  Let E be a normed space. For every subset A of E (resp. B of E'), the polar of A (resp. B) denoted by A° (resp. B°) is the set of \(x' \in E'\) (resp. \(x \in E\) ) for which
\end{enumerate}

\[
\mathcal {R} \left<   x ^ {\prime}, x \right> \geqslant - 1\tag{5}
\]

for all \(x \in A\) (resp. \(x' \in B\) ). When A (resp. B) is a vector subspace, the relation (5) is equivalent to \(\langle x', x \rangle = 0\) , and we then say that \(A^{\circ}\) (resp. \(B^{\circ}\) ) is the orthogonal of A (resp. B).

\begin{enumerate}
\def\labelenumi{\arabic{enumi})}
\setcounter{enumi}{5}
\item
  (« The Bipolar Theorem »). Let E be a normed space. Let A (resp. B) be a subset of E (resp. E') which contains 0. Then the bipolar A°° of A (resp. B°° of B) is the closure for the topology σ(E, E') (resp. σ(E', E)) of the convex envelope of A (resp. B).
\item
  Let A be a subset of a normed space E. Let x be a point in the closure of A with respect to the topology \(\sigma(\mathrm{E}, \mathrm{E}')\) . Then x is the limit (in the norm sense) of a sequence of points of the convex envelope of A. In particular, the convex subsets of E that are closed in the normed space E are the same as those that are closed for \(\sigma(\mathrm{E}, \mathrm{E}')\) .
\item
   Let E be a normed space and M be a vector subspace of E. For every linear form \(u_{0} \in M'\) , there exists a linear form \(u \in E'\) extending \(u_{0}\) and such that \(\|u\| = \|u_{0}\|\) . Let H be the orthogonal of M in \(E'\) ; then the orthogonal \(H^{\circ}\) of H is the closure of M in E.
\end{enumerate}

\paragraph{Transposition}

\begin{enumerate}
\def\labelenumi{\arabic{enumi})}
\setcounter{enumi}{8}
\tightlist
\item
  Let E and F be two normed spaces and \(u \in \mathcal{L}(E; F)\) . The transpose \(^{t}u \in \mathcal{L}(F'; E')\) of u is defined by the relation
\end{enumerate}

\[
\langle u (y ^ {\prime}), x \rangle = \langle y ^ {\prime}, u (x) \rangle \quad \text {   for   all   } x \in E, y ^ {\prime} \in \mathrm{F} ^ {\prime}.\tag{6}
\]

We have \(\|^{t}u\| = \|u\|\) . The kernel of u is the orthogonal in E of the image of \(^{t}u\) . The kernel of \(^{t}u\) is the orthogonal in F' of the image of u.

\begin{enumerate}
\def\labelenumi{\arabic{enumi})}
\setcounter{enumi}{9}
\tightlist
\item
  Let E be a normed space, M be a closed vector subspace of E and \(F = E/M\) . Let i be the canonical injection of M in E and let \(\pi\) be the canonical surjection of E on F. Then \(^{t}i\) has as its kernel the orthogonal \(M^{\circ}\) of M and induces, on passing to the quotient, an isometry of \(E'/M^{\circ}\) on \(M'\) . Further \(^{t}\pi\) is an isometry of \(F'\) on \(M^{\circ}\) .
\end{enumerate}

\paragraph{Conditions for continuity of a linear mapping}

\begin{enumerate}
\def\labelenumi{\arabic{enumi})}
\setcounter{enumi}{10}
\tightlist
\item
  Let E and F be two Banach spaces and u be a linear mapping of E in F. Suppose that for every sequence \((x_{n})_{n\geqslant0}\) of points of E tending to 0 and for which the sequence \((u(x_{n}))_{n\geqslant0}\) has a limit y in F, then y is necessarily 0. Then u is continuous.
\end{enumerate}

* Suppose that for every compact subset K of E, for every positive measure \(\mu\) on K and for every linear continuous form \(y'\) on F, the restriction of \(y' \circ u\) to K is \(\mu\) -measurable. Then u is continuous.*

\begin{enumerate}
\def\labelenumi{\arabic{enumi})}
\setcounter{enumi}{11}
\tightlist
\item
  Let \(\mathbf{E}\) and \(\mathbf{F}\) be two Banach spaces and \(u \in \mathcal{L}(\mathbf{E};\mathbf{F})\) . Then either \(u(\mathbf{E})\) is meagre, or \(u\) is surjective.
\end{enumerate}

Suppose that u is surjective. Then there exists a number C \textgreater{} 0 such that, for all \(y \in F\) , there exists \(x \in E\) with \(u(x) = y\) and \(\|x\| \leqslant C \cdot \|y\|\) . If N is the kernel of u, then u induces on passing to the quotient a homeomorphism of E/N on F.

\begin{enumerate}
\def\labelenumi{\arabic{enumi})}
\setcounter{enumi}{12}
\item
  Let \(\mathbf{E}\) and \(\mathbf{F}\) be two Banach spaces. If \(u\) is a continuous linear mapping of \(\mathbf{E}\) in \(\mathbf{F}\) that is bijective, then \(u^{-1}\) is continuous.
\item
  Let E and F be two Banach spaces, let \(u \in \mathcal{L}(\mathrm{E}; \mathrm{F})\) and \(x' \in E'\) . For \(x'\) to belong to the image of \(^{t}u\) , it is necessary and sufficient that there exists a number C \textgreater{} 0 such that
\end{enumerate}

\[
\left| \langle x ^ {\prime}, x \rangle \right| \leqslant \mathrm{C.} \| u (x) \|\tag{7}
\]

for all \(x \in E\) .

\begin{enumerate}
\def\labelenumi{(\arabic{enumi})}
\setcounter{enumi}{14}
\tightlist
\item
   Let E and F be two Banach spaces and \(u \in \mathcal{L}(E; F)\) . In order that u be surjective, it is necessary and sufficient that there exists a number C \textgreater{} 0 such that \(\|^{t}u(y')\| \geqslant C. \|y'\|\) for all \(y' \in F'\) .
\end{enumerate}

\paragraph{The Banach-Steinhaus Theorem}

\begin{enumerate}
\def\labelenumi{\arabic{enumi})}
\setcounter{enumi}{15}
\item
  (« The Banach-Steinhaus Theorem »). Let E be a Banach space, F a normed space and let \((u_{i})_{i\in\mathbf{I}}\) be a family of elements of \(\mathcal{L}(\mathrm{E};\mathrm{F})\) . Let A be the subset of \(x\in E\) such that \(\sup_{i\in I}\|u_{i}(x)\|<+\infty\) . Then either A is meagre and its complement is dense in E, or alternatively \(\sup_{i\in I}\|u_{i}\|<+\infty\) . In particular, if A = E, then \(\sup_{i\in I}\|u_{i}\|<+\infty\) .
\item
  Let E and F be two Banach spaces and let \((u_{n})_{n\geq0}\) be a sequence of elements of \(\mathcal{L}(\mathrm{E};\mathrm{F})\) . Suppose that the limit \(u(x)=\lim_{n\to\infty}u_{n}(x)\) exists for all \(x\in E\) . Then \(\sup_{n}\|u_{n}\|<+\infty\) , u is continuous and the sequence \((u_{n})\) tends to u uniformly on every compact subset of E.
\end{enumerate}

\paragraph{Properties of the weak topology on a dual}

\begin{enumerate}
\def\labelenumi{\arabic{enumi})}
\setcounter{enumi}{17}
\tightlist
\item
  Let \(\mathbf{E}\) be a Banach space and \(\mathbf{B}'\) be a subset of \(\mathbf{E}'\) . The following conditions are equivalent:
\end{enumerate}

\begin{enumerate}
\def\labelenumi{(\roman{enumi})}
\item
  B' is contained in a ball of E'.
\item
  B' is relatively compact for the topology \(\sigma(\mathrm{E}^{\prime}, \mathrm{E})\) .
\item
  For all \(x \in \mathbf{E}\) , we have \(\sup_{x' \in \mathbf{B}} |\langle x', x \rangle| < +\infty\) .
\end{enumerate}

\begin{enumerate}
\def\labelenumi{\arabic{enumi})}
\setcounter{enumi}{18}
\item
  Let E be a Banach space and let B' be the (closed) unit ball of E'. Then B' is compact for \(\sigma(E', E)\) . Suppose that there exists a countable total subset of E; then B' is metrisable for \(\sigma(E', E)\) , and there exists a countable subset of E' that is dense for \(\sigma(E', E)\) .
\item
  Let \(\mathbf{E}\) be a Banach space, \(u\) be a linear form on \(\mathbf{E}'\) and \(\mathbf{B}'\) be the unit ball of \(\mathbf{E}'\) . The following conditions are equivalent:
\end{enumerate}

\begin{enumerate}
\def\labelenumi{(\roman{enumi})}
\item
  There exists \(\mathrm{x} \in \mathrm{E}\) such that \(u(x') = \langle x', x \rangle\) for all \(x' \in \mathrm{E}'\) .
\item
  The restriction of \(u\) to \(\mathbf{B}'\) is continuous for the topology \(\sigma(\mathrm{E}', \mathrm{E})\) .
\item
  For every sequence \((x_{n}^{\prime})\) of elements of \(E^{\prime}\) that tends to 0 for \(\sigma(\mathrm{E}^{\prime}, \mathrm{E})\) , we have \(\lim_{n \to \infty} u(x_{n}^{\prime}) = 0\) .
\end{enumerate}

\begin{enumerate}
\def\labelenumi{\arabic{enumi})}
\setcounter{enumi}{20}
\tightlist
\item
  Let E be a Banach space, \(B'\) be the unit ball of \(E'\) and C be a convex subset of \(E'\) (in particular a vector subspace). In order that C be closed for \(\sigma(E', E)\) , it is necessary and sufficient that the intersection \(C \cap rB'\) be closed for \(\sigma(E', E)\) for every real number r \textgreater{} 0.
\end{enumerate}

\paragraph{Reflexive spaces}

\begin{enumerate}
\def\labelenumi{\arabic{enumi})}
\setcounter{enumi}{21}
\tightlist
\item
   Let E be a normed space, \(E''\) be its bidual and i be the canonical mapping of E in \(E''\) . The unit ball of \(E''\) is the closure for \(\sigma(\mathrm{E}'', E')\) of the image under i of the unit ball of E.
\end{enumerate}

The following conditions are equivalent :

\begin{enumerate}
\def\labelenumi{(\roman{enumi})}
\item
  The isometric mapping \(i:\mathbf{E}\mapsto \mathbf{E}''\) is surjective.
\item
  The unit ball in E is compact for \(\sigma (\mathrm{E},\mathrm{E}^{\prime})\)
\end{enumerate}

When these conditions are satisfied, we say that E is reflexive.

\paragraph{Topologies compatible with the duality}

\begin{enumerate}
\def\labelenumi{\arabic{enumi})}
\setcounter{enumi}{22}
\tightlist
\item
  Let \(\mathbf{E}\) be a Banach space and \(\mathcal{T}\) a locally convex topology on \(\mathbf{E}\) . The following conditions are equivalent:
\end{enumerate}

\begin{enumerate}
\def\labelenumi{(\roman{enumi})}
\item
  The topology \(\mathcal{T}\) is finer than \(\sigma (\mathrm{E},\mathrm{E}^{\prime})\) and coarser than the topology defined on E by the norm.
\item
  \(\mathbf{E}'\) is the set of linear forms on \(\mathbf{E}\) that are continuous for \(\mathcal{T}\) .
\end{enumerate}

Suppose that these conditions are satisfied. Let A be a subset of E. Then A is relatively compact for T if and only if every sequence of points of A has a cluster point for T in E. If this is so then the balanced convex envelope of A is relatively compact for T.

